% ============================================================
% Capstone Draft 7 — Formatted for submission
% Compile with: pdflatex → pdflatex (run twice for page refs)
% Or upload to Overleaf (overleaf.com) and click Compile.
% ============================================================
\documentclass[twocolumn, 10pt, a4paper]{article}

% ── GEOMETRY ─────────────────────────────────────────────────────────────────
\usepackage[
  top=2.2cm, bottom=2.8cm,
  left=1.6cm, right=1.6cm,
  columnsep=0.75cm
]{geometry}

% ── FONTS AND TYPOGRAPHY ─────────────────────────────────────────────────────
\usepackage{times}      % Times New Roman body text (matches PNAS)
\usepackage{microtype}  % Improved justification / hyphenation

% ── TABLES ───────────────────────────────────────────────────────────────────
\usepackage{booktabs}
\usepackage{array}
\usepackage{tabularx}

% ── LINKS AND CROSS-REFERENCES ───────────────────────────────────────────────
\usepackage[
  colorlinks=true,
  linkcolor=black,
  citecolor=black,
  urlcolor=blue
]{hyperref}
\usepackage{lastpage}   % \pageref{LastPage}

% ── HEADERS / FOOTERS ────────────────────────────────────────────────────────
\usepackage{fancyhdr}

% Fill in your GitHub URL here:
\newcommand{\githublink}{\href{https://github.com/thisisgoingtobechanged/l6_capstone}{github.com/thisisgoingtobechanged/l6\_capstone}}
\newcommand{\LIS}{London Interdisciplinary School}
\newcommand{\thesubdate}{June 2026}
\newcommand{\studentid}{22001041834}
\newcommand{\pagerefs}{\thepage\ of \pageref{LastPage}}

% -- Recurring pages (2 onwards) --
\pagestyle{fancy}
\fancyhf{}
\renewcommand{\headrulewidth}{0pt}
\renewcommand{\footrulewidth}{0.4pt}
\fancyfoot[L]{\footnotesize \githublink\ $|$ ID:~\studentid}
\fancyfoot[R]{\footnotesize \LIS\ $|$ \thesubdate\ $|$ \pagerefs}

% -- Page 1 special --
% NOTE: \thispagestyle must be called AFTER \maketitle (which resets to 'plain').
% It is applied immediately after \twocolumn[...] in the document body.
\fancypagestyle{firstpagestyle}{
  \fancyhf{}
  \renewcommand{\headrulewidth}{0pt}
  \renewcommand{\footrulewidth}{0.4pt}
  \fancyfoot[L]{\footnotesize \githublink}
  \fancyfoot[R]{\footnotesize \LIS\ $|$ \thesubdate\ $|$ \pagerefs}
}

% ── MATHS ────────────────────────────────────────────────────────────────────
\usepackage{amsmath}
\usepackage{amssymb}  % \blacktriangleleft for citation direction notation

% ── MISC ─────────────────────────────────────────────────────────────────────
\usepackage{parskip}    % Space between paragraphs instead of indent
\setlength{\parskip}{4pt}
\usepackage{enumitem}
\setlist[itemize]{leftmargin=*, topsep=2pt, itemsep=1pt}

% ── TITLE ────────────────────────────────────────────────────────────────────
\title{\textbf{The World Leader Effect in Animal Welfare Policy:}\\
  \large Categorising evidence of diffusion mechanisms in primary animal welfare
  legislation in the Danish, Austrian and German parliaments, 2009--2019}
\author{Student ID: \studentid \quad \LIS}
\date{\thesubdate}

% ============================================================
\begin{document}
% ============================================================

% Title and preamble span both columns
\twocolumn[{%
\begin{@twocolumnfalse}

\maketitle
\vspace{-0.5em}
\hrule
\vspace{0.8em}

% ── ABBREVIATIONS + COUNTRY CODES ────────────────────────────────────────────
\begin{minipage}[t]{0.45\textwidth}
\textbf{Abbreviations}\\[4pt]
\begin{tabular}{ll}
\toprule
Acronym & Meaning \\
\midrule
FAW & Farmed Animal Welfare \\
EU  & European Union \\
\bottomrule
\end{tabular}
\end{minipage}%
\hfill
\begin{minipage}[t]{0.45\textwidth}
\textbf{Country codes (ISO 3166)}\\[4pt]
\begin{tabular}{ll}
\toprule
Code & Country \\
\midrule
DK & Denmark \\
AT & Austria \\
DE & Germany \\
NL & Netherlands \\
SE & Sweden \\
CH & Switzerland \\
UK & United Kingdom \\
\bottomrule
\end{tabular}
\end{minipage}

\vspace{0.6em}
\noindent\small\textit{In this study, `A $\blacktriangleleft$ B' indicates country A is citing country B in its legislative process.}
\vspace{0.6em}
\hrule
\vspace{0.8em}

% ── ABSTRACT ─────────────────────────────────────────────────────────────────
\noindent\textbf{Abstract} \quad
This pilot study asks how three Northern European countries make reference to each
other and world leaders in their primary legislation\footnote{Primary legislation is
law created by the legislative branch of government. In parliamentary democracies,
this will generally consist of parliaments passing statutes or acts that lay out
general principles. From this, the executive branch then issues secondary
legislation, which is more specific and legally enforceable.} relating to farmed
animal welfare between 2009 and 2019. Using the ParlLawSpeech V2 parliamentary
corpus (Schwalbach et al., 2025), 453 documents, consisting of bills and laws, were
assembled from the legislature of Germany, Austria, and Denmark. A
regular-expression citation extraction pipeline produced 355 foreign country
reference contexts, of which 62 formed the analytical sample after AI-assisted
classification using the taxonomy of De Ruiter and Schalk (2017). Germany emerges
as a discursive leader among follower countries, receiving 22 within-trio analytical
citations, 20 of which come from Denmark. Qualitative analysis reveals that
citations are mostly informational across all country pairs. This indicates that the
operative diffusion mechanism for farmed animal welfare in the EU is learning rather
than emulation. This is despite Denmark adopting most welfare reforms at the same
time as or before Germany. Discursive attention that indicates learning under
De~Ruiter and Schalk's taxonomy may therefore also or instead be tracking competitive
economic dependency. Future research should further this programme of improving the
measures by which the diffusion mechanism of animal welfare policy can be unpicked as
existing measures are still confounding.

\vspace{1.2em}
\end{@twocolumnfalse}
}]
\thispagestyle{firstpagestyle}  % override \maketitle's reset to 'plain'

% ============================================================
% TWO-COLUMN BODY STARTS HERE
% ============================================================

\section{Background}

Over the last few decades, consumer concern regarding the implications of industrial
animal agriculture has risen (Broom, 2017, p.~18; European Court of Auditors, 2018,
p.~3). The EU moved quickly in the late 1990s and early 2000s to introduce minimum
welfare standards (Broom, 2017, p.~24), and recognised animals as `sentient beings'
in the 2009 Treaty of Lisbon (European Court of Auditors, 2018, p.~3; Broom, 2017,
p.~26). Supranational EU legislation has since stalled. Dairy cows, farmed fish, and
rabbits are unprotected by specific EU directives (Broom, 2017, p.~54). European
Citizens' Initiatives for a `Fur Free Europe' and to `End the Cage Age' both
succeeded in gaining the requisite signatures but the European Commission is yet to
bring forward a legislative proposal on either (European Commission, 2021, 2023).

Even the least progressive EU states remain comparatively advanced in the global
context because of the EU setting a floor for animal welfare standards. However, EU
minimums now often lag behind the progressive legislation of some of its member
states and there is considerable variation between leader countries and laggards even
within the EU (Vogeler, 2017). Five out of the six highest rated countries worldwide
(ranked according to the World Animal Protection Index (API), each graded with a `B'
overall) are current or former EU member states (World Animal Protection, 2020).
These are Austria, Denmark, Sweden, the Netherlands and the UK. The sixth, not an EU
member, is Switzerland. It is the only country to have received the highest `A' grade
from the API, in 2014\footnote{Switzerland was demoted to a `B' grade in 2020, the
same rank as this study's other World Leaders in animal welfare.} (World Animal
Protection, 2020). These six countries are this study's World Leaders.

\section{Literature Review}

Policy diffusion is `the process whereby policy choices in one unit are influenced
by policy choices in other units' (Maggetti \& Gilardi, 2016, p.~87). There is
consensus in the literature that diffusion is a product of interdependence. However,
this interdependence can take different forms (Maggetti \& Gilardi, 2016, p.~87).
The literature refers to these forms as `mechanisms', which in over a thousand
articles (Shipan \& Volden, 2012, p.~1) have been grouped into four broad
categories: learning, emulation, competition and coercion (Shipan \& Volden, 2008,
p.~840), which for clarity are referred collectively to in this study as the LECC
model.

Learning is defined as where `the policies in one unit are influenced by the
consequences of similar policies in other units' (Maggetti \& Gilardi, 2016,
p.~4). Emulation by contrast is not related to the objective consequences of a
policy and instead its symbolic and socially-constructed characteristics are crucial
(Maggetti \& Gilardi, 2016, p.~5). Competition is the reaction between units in an
attempt to attract and retain resources (Maggetti \& Gilardi, 2016, p.~5). An
imagined example of each might be: `Swedish policy has proven effective in reducing
sow mortality' (learning) versus `we should be more like Sweden' (emulation) versus
`we need to mimic Sweden before Swedish producers undercut us' (competition). Although
coercion is not a mechanism this study examines, it is defined as when threats are
made to subordinate units if they do not comply with policy changes suggested by
the senior unit, e.g.\ `The EU will fine us if we do not adopt animal welfare
standards'.

Graham, Shipan and Volden (2013) conducted a systematic review documenting how
diffusion research has itself `diffused' across the social sciences by means of the
shared conceptual vocabulary of these four mechanisms. Mallinson (2021) extended this
in a meta-review of policy diffusion studies in American states, estimating effect
sizes of the most commonly used diffusion variables and identifying biases in
published research. The cumulative picture is that diffusion effects are real and
robust, but their magnitude and operative mechanisms vary substantially across policy
types and political contexts. Maggetti and Gilardi (2016) provide a meta-analysis of
114 studies demonstrating that the same indicators are used to measure different
diffusion mechanisms and different indicators are used to measure the same diffusion
mechanism (Maggetti \& Gilardi, 2016, p.~18).

There are multiple criticisms of the LECC model in the literature which have resulted
in alternative proposals to redress its perceived shortcomings. The first, the
underspecification of these four terms, has led to the development of new definitions
to be applied in a more strictly deductive fashion. Blatter et al.\ (2021) address
bias, and reformulate the mechanisms into interest-driven, rights-driven,
ideology-driven, and recognition-driven pathways. Gilardi and Wasserfallen (2019) are
concerned by a failure of omission by the model. They argue that the LECC model
underspecifies the political dimensions of the process, and treats the selective
borrowing of foreign policy as a well-calculated choice by legislators trying to
garner legitimacy for their proposals. Genovese et al.\ (2016) corroborate this and
show that such an assumption would be a systematic error. Governments are rarely
presented with a single diffusing policy in isolation. There are frequently
alternative policies to consider, but a country must also consider their dependency on
foreign resources and will choose a substitute policy by evaluating the more limited
options available under these constraints.

Linos argues instead that diffusion and domestic democracy are mutually reinforcing
(Linos, 2013, p.~2). Politicians emulate the policy choices of familiar, wealthy and
culturally proximate countries to signal to `uninformed' voters that their proposals
are mainstream, competently designed and vetted (Linos, 2013, p.~3).

Despite these legitimate criticisms and modifications of the most widespread
definition framework (LECC), this study chooses to persist with using the LECC
definitions because in so doing `we open our work up to more natural comparisons to
similar studies elsewhere' (Graham et al., 2013, p.~700). Placing this capstone in a
lineage of studies that share these familiar terms offers advantages. It can build on
existing insights as well as enabling more straightforward verification and
reproduction of results. To mitigate the pitfalls of the LECC framework, this study
requires explicit and systematic reflection between how concepts are being
conceptualised and how they will be operationalised. This will be laid out in the
methodology. The challenge of the LECC framework is in how to consistently test its
mechanisms in different contexts. The background definitions of these concepts, before
they are systematised, are uncontentious.

Policy diffusion operates across multiple governance levels. Countries can diffuse
policies between one another in a peer-to-peer horizontal process (De Ruiter \&
Schalk, 2017, p.~2). Particularly pertinent in the European Union context is vertical
influence. `Bottom-up' diffusion occurs when individual member states influence the EU
parliamentary process. `Top-down' diffusion involves EU mandates reshaping member
states' behaviour (Shipan \& Volden, 2006, p.~825). It is in top-down instances that
coercion is most commonly observed. Now that EU-level farmed animal welfare
legislation has stalled, the operative diffusion mechanism has therefore shifted from
top-down to peer-to-peer.

Policy diffusion literature leverages the methods of analysis of many different social
science disciplines. All the mechanisms described above have social, political and
economic dimensions to them. The significant influence of these contextual factors
makes it difficult to develop accurate, predictable theories of policy diffusion at
scale. It is therefore essential to stratify the analysis as much as possible by
limiting the research parameters to well-established geographic or federal blocs e.g.\
The USA or the EU. The focus is less on the specific policy or category under
investigation and more about developing a better theory of the \textit{mechanics} of
policy diffusion between similar countries or states. Interestingly, although this
recognition is well-aligned to the axioms of complexity science and its pursuit of
meta-parsimonious explanations, the author did not find any literature on policy
diffusion or animal welfare policy conducted through that prism.

The terms `leader' and `follower' are widely used in policy diffusion literature. In
different instances, a policy leader is defined as: the first mover in an adoption
sequence; the most cited or emulated country in the debates of adopting states
(discursive leadership); the state whose legislative turn of phrase is most
extensively reused; or the state exerting the greatest causal influence on the
probability of adoption elsewhere. This study focuses on the second definition of a
policy `leader' as the country most frequently cited when legislators in other
jurisdictions debate animal welfare legislation\footnote{This is for reasons expanded
on in the methodology, but briefly because it can be measured directly from
parliamentary records using the citation analysis methods described in the
methodology section.}.

Members of Parliament operate in an environment of `attention scarcity' (Jones \&
Baumgartner, 2005, as cited in De Ruiter \& Schalk, 2017, p.~1), meaning they must
find ways to lower `information collection costs', in other words reduce complexity
when identifying the most critical issues to address (De Ruiter \& Schalk, 2017,
p.~1). One key strategy is referencing the policy practices of other countries during
plenary debates. De Ruiter and Schalk (2017, p.~4) demonstrate that MPs use new
information about the policy practice in another country in parliamentary debate as
it becomes available and citations could be a lever of legitimation. Linos (2013,
p.~3) provides a theoretical basis for this, hypothesising that this is to signal a
proposal is mainstream and internationally vetted, and therefore normatively
acceptable to peers (and voters). There are other sources of policy that could be
used as exemplars (e.g.\ think-tanks) but they are perhaps perceived as less
rhetorically effective in debate.

De Ruiter and Schalk (2017) inductively develop a taxonomy of foreign country
references in Dutch plenary speeches that distinguishes between informational
citations (where a foreign policy is described to inform deliberation) and
legitimating citations (where the foreign case is invoked to validate a domestic
proposal). A third type, cautionary citations, cite a foreign case as a warning or
negative example. Note that citation frequency analysis captures salience in the
debate, not necessarily cause and effect. A country can be cited frequently as a
cautionary tale. It can be cited because it is culturally proximate and well-known
rather than because its specific policy design has been substantively vetted and
appropriated (i.e.\ whether the policy is diffusing via emulation or learning).

Policy diffusion does not solely occur at the final stage of policy adoption. A
significant but often neglected aspect of the diffusion process is the issue-definition
stage (Gilardi et al., 2021, p.~21). Before policies are adopted or placed on an
agenda, issues must be identified and defined through specific policy frames (Gilardi
et al., 2021, p.~22). Studying diffusion solely through the lens of successful
adoptions leads to a `pro-innovation bias'. Policies pass through various framing
stages and adoptions themselves are rare (Gilardi et al., 2021, pp.~22--23). This
study, therefore, includes proposed bills that did not pass alongside enacted laws.

This research focuses specifically on laws and bills concerning farmed animal welfare
in Denmark, Germany, and Austria between 2009 and 2019. This ten-year period is the
overlap between the available data from ParlLawSpeech for these three countries.

\noindent From this literature review the following sub-questions have emerged:

\begin{enumerate}[leftmargin=*, topsep=2pt, itemsep=1pt]
  \item Which foreign countries appear most frequently in the animal welfare
        legislation of each citing country, and does this pattern differ across the
        trio?
  \item Do those citations function primarily as evidence-gathering (informational),
        social proof (legitimating), or competitive awareness (cautionary)?
  \item Does the direction of discursive attention align with the sequence in which
        farmed animal welfare legislation was adopted?
\end{enumerate}

\section{Methodology}

This study uses a convergent parallel mixed-methods design (Creswell \& Plano Clark,
2011), in which a quantitative strand (citation frequency) and a qualitative strand
(citation classification) are run in parallel and integrated in a final synthesis
table.

ParlLawSpeech is a multilingual parliamentary corpus covering legislative documents
from multiple European countries (Rauh et al., 2023). For this study, the relevant
components are the bills and enacted laws from Germany, Austria, and Denmark. The
corpus provides documents in their original languages along with metadata including
adoption date, initiation date and an identifier for each legislative item. This last
links all documents belonging to the same legislative process. It can be used to
recover laws whose keyword coverage in the enacted text is lower than in the
originating bill.

Foreign country citations were extracted using regular expression matching on the
full text of each document in the corpus. Country terms were defined in a
\texttt{COUNTRY\_NAMES} dictionary covering each target country's name in German,
Danish, and English, with word-boundary anchors (\texttt{\textbackslash b}) to
prevent false matches. This prevents compound words such as
\textit{deutschnationale}, [German nationalist] from registering as a Germany
reference. The target countries were Germany, Austria, Denmark, Sweden, Switzerland,
the Netherlands, and the United Kingdom. Sweden and Switzerland were included as
reference countries known from the adoption sequencing literature (Wallenbeck et al.,
2024) to be relevant to the regional welfare diffusion context.

For each regex match, a context window of $\pm600$ characters around the match was
extracted to enable subsequent classification. Each match generated one row in the
citation dataset: one document can contribute multiple rows if the target country is
mentioned more than once.

Each of the 355 extracted citation contexts was classified using the Claude Haiku 4.5
language model via the Anthropic API. The classification scheme was based on De~Ruiter
and Schalk's (2017) taxonomy for plenary speeches, in this case deployed on bills and
laws and extended to include three exclusion categories.

\noindent\textbf{Analytical categories (can co-occur).}
\begin{itemize}
  \item \textit{Informational:} the foreign case is cited to provide factual
        information about a policy elsewhere; function is to inform deliberation.
  \item \textit{Legitimating:} the foreign case is cited to validate or authorise a
        domestic proposal; function is to signal the proposal is mainstream and
        peer-vetted.
  \item \textit{Cautionary:} the foreign case is cited as a warning or negative
        example.
\end{itemize}

\noindent\textbf{Exclusion categories (mutually exclusive).}
\begin{itemize}
  \item \textit{False positive:} the country name appears but not as a reference to
        that country as a political actor or policy exemplar (e.g.\ language
        references such as \textit{deutsche \"Ubersetzung}, proper nouns, geographic
        descriptions).
  \item \textit{Out of scope:} the country is cited as a policy actor, but the
        document concerns non-farmed animal welfare (companion animals, zoo animals,
        wildlife).
  \item \textit{EU implementation:} the document is transposing or ratifying EU
        legislation; the country is mentioned as a co-signatory, not as a voluntary
        peer whose policy is being studied.
\end{itemize}

Before examining the AI output, the author independently classified 15 randomly
sampled contexts and recorded labels. The sample was locked.

\section{Results}

The filtering pipeline retained 453 documents in total: 204 laws and 249 bills.
Denmark contributes the largest share of both categories (laws: 56\%, bills: 51\%).
This reflects the larger size of the Danish component of ParlLawSpeech for this
period. Germany and Austria each contribute roughly equally to the remainder.

Regular expression extraction produced 355 foreign country citation contexts across
the 453 documents. Denmark generated the largest share of raw citations (213, 60\%),
followed by Austria (91, 26\%) and Germany (51, 14\%). This asymmetry partly reflects
Denmark's larger corpus share and the denser text of Danish bills. Denmark produces
roughly twice as many foreign country mentions per document as Germany (0.89 versus
0.43 citations per document). Austria's per-document rate (0.96) is slightly higher,
though Austria's smaller corpus means its absolute count remains far behind Denmark's.

The most-cited country by raw frequency is Switzerland (94 citations), followed by
Germany (93) and Sweden (66). Austria receives 22 raw citations; Denmark receives 10;
the Netherlands 33; the UK 37.

Switzerland's high raw count reflects EU implementation documents. Switzerland is
frequently named as a co-signatory or comparative case in EU documents. This is
consistent with Switzerland's institutional position vis-\`{a}-vis the EU. Since
1995, Switzerland has pursued a policy of autonomous regulatory alignment with EU
veterinary and food safety law, formalised from 2002 in the bilateral agreement on
trade in agricultural products and extended in 2009 to a Common Veterinary Area
covering animal health and welfare standards.

Of the 355 citation contexts, 293 (82.5\%) were excluded:
\begin{itemize}
  \item \textit{EU implementation:} 87 (24.5\%)
  \item \textit{False positive:} 102 (28.7\%)
  \item \textit{Out of scope:} 104 (29.3\%)
\end{itemize}

The remaining 62 contexts (17.5\%) formed the analytical sample. The high exclusion
rate is expected. Country names appear frequently in legislative documents for
incidental reasons, such as a part of EU regulatory language or in legislation about
zoo or companion animal welfare.

The distribution of labels across the 62 analytical citations is:
\begin{itemize}
  \item \textit{Informational:} 58
  \item \textit{Cautionary:} 6
  \item \textit{Legitimating:} 2
\end{itemize}

Analytical categories can co-occur (and do, in this case). Therefore there are more
than 62 citations.

Across all analytical citations in the corpus, the dominant function of foreign
country references is informational. Using De Ruiter and Schalk's (2017) framework
it appears that citations of other countries are primarily to describe what those
countries have already done.

\begin{table}[h]
\centering
\footnotesize
\caption{Citation mechanisms by country pair}
\label{tab:mechanisms}
\begin{tabular}{lrrrrl}
\toprule
Citing $\blacktriangleleft$ Cited & N & Inf.\% & Leg.\% & Caut.\% & Mechanism \\
\midrule
DK $\blacktriangleleft$ DE & 20 & 90 & 0 & 20 & Learning \\
DE $\blacktriangleleft$ DK & 3  & 100 & 0 & 0 & Learning \\
AT $\blacktriangleleft$ DE & 2  & 100 & 50 & 0 & Learning + Emulation \\
DE $\blacktriangleleft$ AT & 2  & 50 & 50 & 50 & Learning + Emulation \\
AT $\blacktriangleleft$ DK & 1  & 100 & 0 & 0 & Learning \\
DK $\blacktriangleleft$ AT & 0  & --- & --- & --- & --- \\
\bottomrule
\end{tabular}
\end{table}

Germany receives 22 analytical citations, 20 from Denmark and 2 from Austria.
Denmark receives 4, 3 from Germany and 1 from Austria. Austria receives 2, both from
Germany.

Denmark continues to cite higher than Austria and Germany when extended to World
Leaders. The full integration table records DK $\blacktriangleleft$ SE (11 analytical citations) as
the second-strongest citation relationship in the dataset. Denmark's outward
orientation is not exclusively directed at Germany. Sweden's prominence in Danish
legislative discourse is consistent with its role as the earliest adopter of pig
welfare reforms in the region.

Across all analytical citations ($N = 62$), Denmark accounts for 40 of them (65\%).
Germany 14, and Austria 8. Denmark is the most prolific citing country and the
recipient of the fewest within-trio citations. Germany is the least prolific citing
country and the most-cited.

The raw Cohen's Kappa, comparing the author's labels to the AI's classifications,
was $\kappa = 0.281$. This suggests fair agreement on the Landis and Koch (1977)
scale.

This figure is misleading in a specific and important way. Eleven of the fifteen
validation cases concerned the boundary between \textit{false positive} and
\textit{out of scope}. These categories both lead to exclusion from the retained
sample. The author classified several cases as \textit{out of scope} that the AI
classified as \textit{false positive}, or vice versa. This boundary is admittedly
ambiguous. This is because a document about zoo animal welfare that mentions Germany
could be argued either way, depending on whether the country is read as a genuine
policy exemplar or merely an incidental mention.

When the three exclusion categories are collapsed into a single `excluded' label,
agreement on the analytically consequential include-versus-exclude decision rises to
14 out of 15 cases (93.3\%). The Cohen's Kappa on this binary include/exclude
decision is $\kappa = 0.762$\footnote{%
$0.762 = (0.933 - P_e) / (1 - P_e) \therefore P_e \approx 0.719$}. The author
classified one case as \textit{unclear} (which would have included it in the
analytical sample) while the AI classified it as \textit{false positive}, which
excluded it.

The two cases that both author and AI agreed should be included in the analytical
sample were independently assigned the `informational' label. One case where the
author assigned `cautionary' in addition to informational was classified by the AI
as informational only.

These results were checked against temporal adoption sequencing data from Wallenbeck
et al.\ (2024), who documented the year in which each country adopted specific pig
welfare measures relative to Sweden. For the study period, the relevant in-corpus
measures are gestation housing (Austria in 2012, Germany in Denmark in 2013). This
was the only welfare adoption event occurring for all three study countries within
the 2009--2019 window at a point captured by the corpus.

\subsection{Diffusion Mechanisms}

The dominant mechanism across the corpus is learning according to the deductive
framework. 93.5\% of analytical citations are informational, whereas legitimating
citations make up 3.2\% (2/62). The most extensively cited relationship (Denmark to
Germany, $N = 20$) is 90\% informational, 0\% legitimating, and 10\% cautionary.

All six cautionary classifications identified overall appear in three specific
citation relationships: DK $\blacktriangleleft$ DE, DE $\blacktriangleleft$ AT, and DE $\blacktriangleleft$ NL. The DK $\blacktriangleleft$ DE
cautionary cases concern disease outbreak risk. Danish legislation cites Germany's
2011--2012 dioxin crisis and historical swine fever outbreaks as exemplars of what
can happen when food safety oversight fails. The DE $\blacktriangleleft$ AT and DE $\blacktriangleleft$ NL
cautionary cases concern fur farming. A 2017 German bill cites Austria, the
Netherlands and others as countries that have already banned mink farming, positioning
Germany as a laggard relative to its European peers.

Switzerland generates the highest number of raw citations across the corpus (94) but
only 6 analytical citations after exclusion. Germany also cites Switzerland ($N = 3$)
and the Netherlands ($N = 3$) in its analytical sample, primarily in the fur farming
context. A 2017 German bill on mink welfare (mentioned in Vignette~3, below)
compares German minimum space requirements unfavourably with Swiss, Austrian, and
Dutch provisions. This document is responsible for the entirety of the DE $\blacktriangleleft$ AT
and DE $\blacktriangleleft$ NL cautionary classifications.

\subsection{Illustrative Vignettes}

\noindent\textbf{Vignette 1: Cautionary, DK $\blacktriangleleft$ DE (2012, disease control)}

A 2012 Danish bill on veterinary food safety controls invokes Germany's dioxin
contamination crisis to argue for more robust domestic oversight. The context reads:
\textit{`det danske kontrolsystem hidtil vist sig s\aa{} effektivt, at vi har
und\aa{}et kemiske forureninger'}. The Danish control system has so far proven so
effective that Denmark has avoided the chemical contamination incidents [of the kind
seen in Germany]. Germany is cited as a warning. Denmark defines its own system's
quality by contrast with Germany's failure. This is a competitive-awareness citation
with a strong cautionary function. The AI classification (confidence: high)
characterises this as a case where Germany's experience is invoked `to warn against
potential economic consequences Denmark might face if similar contamination occurs,
establishing a cautionary comparative reference.'

\noindent\textbf{Vignette 2: Informational, DK $\blacktriangleleft$ DE (2018, veterinary medicine)}

A 2018 Danish veterinary medicine bill surveys how comparable EU countries handle
the dispensing of measured quantities of veterinary drugs. The relevant passage names
\textit{`Belgien, Finland, Tyskland og Holland'}, Belgium, Finland, Germany, and the
Netherlands, as countries where such dispensing schemes exist in some form. It also
notes that in all these countries (as in Denmark) vets can already administer precise
doses during treatment. Germany appears here as one data point in a comparative
audit, cited descriptively and without evaluative language. The AI classification
(confidence: high) characterises this as `providing factual information about how
other member states handle medication portioning without invoking Germany as a model
for legitimation or warning.'

\noindent\textbf{Vignette 3: Cautionary, DE $\blacktriangleleft$ AT (2017, fur farming)}

A 2017 German bill on fur animal welfare cites Austria, the Netherlands, the United
Kingdom, and several other European countries as jurisdictions that have already
enacted fur farming bans. The context notes that these countries have moved ahead on
animal protection standards. Germany however, despite recommendations to the contrary
from mammalogy experts, still requires only one square metre of space per mink.
Austrian and Swiss requirements are many times larger. Austria is cited as evidence
that Germany is falling behind. The cautionary function is self-directed. The AI
classification (confidence: high) reads this as functioning `both to inform about
existing policy elsewhere and to warn that Germany is lagging behind peer nations in
animal protection standards.'

\section{Discussion}

The results are consistent with Mallinson's (2020) finding that ideological learning
becomes more prominent in later stages of `the [policy] life course', displacing the
regional proximity effects dominant in early stages.

Germany receives the most within-trio citations in farmed animal welfare legislation,
driven overwhelmingly by Denmark. On the standard operational definition used in this
study, Germany is the discursive leader.

However, Germany does consistently not lead Denmark in welfare adoption. Denmark
adopted tail docking and manipulable material restrictions earlier than Germany and
reached gestation housing reform at the same time. There is no empirical basis for
treating Denmark as a discursive follower of a policy pioneer, yet Denmark cites
Germany more than Germany cites Denmark, and more than either Austria or Germany cite
each other.

One plausible explanation is competitive monitoring. Germany is Denmark's largest
neighbour, a major pork-producing economy, and a competitor in European export
markets (Denver et al., 2021). Danish legislators have incentives to track German
welfare legislation to know when Germany will be subject to the same costs that
Danish producers already bear. If Danish pigs are raised under stricter welfare
conditions than German pigs, Danish pork is more expensive to produce (Vogeler,
2017). The moment Germany adopts equivalent standards, that cost disadvantage may
narrow. Denmark's parliament has an interest in knowing when Germany will reach that
point.

This pattern is consistent with recent comparative work on regulatory competition in
farm animal welfare elsewhere. Z\"{o}llmer (2025), examining the diffusion of welfare
standards across US states, argues against the assumption that competitive exposure
straightforwardly suppresses unilateral welfare improvement. The governance mode
through which a standard is implemented determines whether competitive pressure
produces resistance or accommodation (\textit{cf.}~Vogel's (1995) account of `races
to the top' in environmental regulation).

Legitimating citations across the corpus are extremely rare, while De Ruiter and
Schalk (2017) found sufficient legitimating citations to inductively invent a
category for them in the Dutch parliamentary context. The author offers several
possible explanations. The sample is made up by bills and laws, not speeches.
Legitimating rhetoric may be more characteristic of oral parliamentary debate than of
written legislative texts, where the primary function is regulatory specification. A
study of parliamentary speeches may reveal more legitimating citation. Another is
that the countries in this study are still relatively high-welfare both in the
European and global contexts. Legitimating citations may be more characteristic of
laggard countries catching up. A study that included lower-welfare EU member states
as citing countries would enable a direct test of this life-course interpretation.

\subsection{Limitations}

Several limitations should be noted. First, the corpus excludes parliamentary
speeches for the primary analysis, meaning that the full discursive record is not
examined. Second, the Kappa of 0.281 reflects genuine disagreement between the
author and the AI on exclusion subcategory boundaries, even if the include/exclude
agreement was 93.3\%. The AI may have overclassified or underclassified some
analytical cases in ways that affected the mechanism distribution. Third, the small
$N$ in some country pairs (AT $\blacktriangleleft$ DK: 1, AT $\blacktriangleleft$ DE: 2, DE $\blacktriangleleft$ AT: 2) means
that mechanism classifications for those pairs are indicative rather than robust.
Fourth, some DK $\blacktriangleleft$ SE citations in the analytical sample appeared on closer
inspection to concern non-welfare topics (aviation regulation, excise duties) present
in multi-topic bills that passed the keyword filter; this suggests the broad-tier
keyword filter includes some documents with only marginal welfare relevance, and that
the out-of-scope classifier did not catch all instances.

Fifth, the study's chosen taxonomy was developed inductively by De Ruiter and Schalk
(2017) from Dutch parliamentary debate on education and environmental policy. The
categories of informational, legitimating, and cautionary citation may carry
different rhetorical weight in German, Danish and Austrian legislative systems, and
in the specific domain of farm animal welfare. The near-absence of legitimating
citation in this corpus could also reflect a difference between nations in how
foreign references function. This study cannot disentangle this. A validation of the
taxonomy against a corpus from a country or policy domain where De Ruiter and
Schalk's original categories were directly observed would help establish how portable
the taxonomy is.

Future research should extend the analysis to parliamentary speeches, which would
provide a larger and more rhetoric-rich corpus. Including lower-welfare EU member
states as citing countries would enable comparison of legitimating citation rates
across welfare levels. It would provide a more direct test of the life-course
interpretation offered in the discussion. A systematic review of the legislative
language borrowed across jurisdictions would complement the citation frequency
approach used here and triangulate a clearer measure of policy transfer.

\section{Conclusion}

Using the ParlLawSpeech corpus, 453 legislative documents were assembled and 355
foreign country citation contexts were extracted and classified using an AI-assisted
taxonomy derived from De Ruiter and Schalk (2017). Germany emerges as the discursive
leader, receiving 22 within-trio analytical citations compared to 4 for Denmark and
2 for Austria. The mechanism appears to be predominantly learning, as informational
citation accounts for 93.5\% of analytical cases, with legitimating citations very
rare. The cautionary cases concern two distinct phenomena: disease outbreak risk
(DK $\blacktriangleleft$ DE) and laggard positioning relative to peers in fur farming welfare
standards (DE $\blacktriangleleft$ AT, DE $\blacktriangleleft$ NL).

However, the study demonstrates an asymmetry between adoption sequencing and what
discursive attention would indicate under De Ruiter and Schalk's taxonomy. Denmark
adopted several welfare measures at the same time as or before Germany, yet Denmark
monitors Germany's legislative debates far more intensively than Germany monitors
Denmark's. This suggests that discursive leadership, as measured by citation
frequency, may reflect competitive economic dependency, and not learning. Germany may
function as Denmark's discursive reference point due to this economic dependency
rather than because German welfare policy is more advanced.

Countries cite each other for different reasons at different stages of the diffusion
lifecycle. A country that cites frequently may be learning, emulating, or watching a
competitor. Disambiguating these mechanisms is essential for connecting the
legislative process to how policy diffusion occurs.

\vspace{0.5em}
\noindent\small\textit{This work is 5{,}905 words in length from the title to the end of the reference list.}

\vspace{0.3em}
\noindent\small\textit{Verification data, code, supplementary information, and the Capstone product are available at \githublink.}

% ============================================================
% REFERENCES
% ============================================================
\section*{References}
\begingroup
\setlength{\parskip}{3pt}
\footnotesize

\noindent Blatter, J., Schmid, S., \& Bl\"attler, A.~C. (2021). Conceptualising and
measuring political diffusion: A new approach. \textit{Political Studies Review},
\textit{19}(1), 79--95.

\noindent Broom, D.~M. (2017). \textit{Animal welfare in the European Union}.
European Parliament Policy Department for Citizens' Rights and Constitutional
Affairs.

\noindent Creswell, J.~W., \& Plano Clark, V.~L. (2011). \textit{Designing and
conducting mixed methods research} (2nd ed.). SAGE.

\noindent De Ruiter, R., \& Schalk, J. (2017). Cross-national citation in European
national parliaments: A case study of foreign references in Dutch parliamentary
debate. \textit{Acta Politica}, \textit{52}(1), 1--23.

\noindent Denver, S., Christensen, T., Nordstr\"{o}m, J., Lund, T.~B., \&
Sand\o{}e, P. (2021). Is there a potential international market for Danish welfare
pork? A consumer survey from Denmark, Sweden, and Germany. \textit{Meat Science},
\textit{182}, 108633.

\noindent European Commission. (2021). \textit{Questions and answers: Commission's
response to the European Citizens' Initiative on ``End the Cage Age''} [Press
release]. European Commission Press Corner.

\noindent European Commission. (2023). \textit{Commission reply to the `Fur Free
Europe' European Citizens' Initiative} [Communication]. European Citizens'
Initiative Portal.

\noindent European Court of Auditors. (2018). \textit{Animal welfare in the EU:
Closing the gap between ambitious goals and practical implementation} (Special Report
31/2018). Publications Office of the European Union.

\noindent Genovese, F., Schneider, G., \& Wasserfallen, F. (2016). The politics of
EU differentiated integration. \textit{British Journal of Political Science},
\textit{46}(1), 1--22.

\noindent Gilardi, F., Shipan, C.~R., \& W\"{u}est, B. (2021). Policy diffusion:
The issue-definition stage. \textit{American Journal of Political Science},
\textit{65}(1), 21--35.

\noindent Gilardi, F., \& Wasserfallen, F. (2019). The politics of policy diffusion.
\textit{European Journal of Political Research}, \textit{58}(4), 1245--1256.

\noindent Graham, E.~R., Shipan, C.~R., \& Volden, C. (2013). The diffusion of
policy diffusion research in political science. \textit{British Journal of Political
Science}, \textit{43}(3), 673--701.

\noindent Landis, J.~R., \& Koch, G.~G. (1977). The measurement of observer
agreement for categorical data. \textit{Biometrics}, \textit{33}(1), 159--174.

\noindent Linos, K. (2013). \textit{The democratic foundations of policy diffusion:
How health, family, and employment laws spread across countries}. Oxford University
Press.

\noindent Maggetti, M., \& Gilardi, F. (2016). Problems (and solutions) in the
measurement of policy diffusion mechanisms. \textit{Journal of Public Policy},
\textit{36}(1), 87--107.

\noindent Mallinson, D.~J. (2020). Policy innovation adoption across the diffusion
life course. \textit{Policy Studies Journal}, \textit{49}(2), 1--25.

\noindent Mallinson, D.~J. (2021). Growth and gaps: A meta-review of policy
diffusion studies in the American states. \textit{Policy \& Politics},
\textit{49}(3), 369--389.

\noindent Rauh, C., Bock, O., Ceron, A., Froio, C., \& Zhelyazkova, A. (2023).
ParlLawSpeech: A new comparative dataset on parliamentary speeches and legislative
documents. \textit{Legislative Studies Quarterly}, \textit{48}(2), 411--433.

\noindent Schwalbach, J., Hetzer, L., Proksch, S.-O., Rauh, C., \& Seb\H{o}k, M.
(2025). \textit{ParlLawSpeech} [Dataset]. GESIS, Cologne.
\url{https://doi.org/10.7802/2824}

\noindent Shipan, C.~R., \& Volden, C. (2006). Bottom-up federalism: The diffusion
of antismoking policies from U.S.\ cities to states. \textit{American Journal of
Political Science}, \textit{50}(4), 825--843.

\noindent Shipan, C.~R., \& Volden, C. (2008). The mechanisms of policy diffusion.
\textit{American Journal of Political Science}, \textit{52}(4), 840--857.

\noindent Shipan, C.~R., \& Volden, C. (2012). Policy diffusion: Seven lessons for
scholars and practitioners. \textit{Public Administration Review}, \textit{72}(6),
788--796.

\noindent Vogel, D. (1995). \textit{Trading up: Consumer and environmental regulation
in a global economy}. Harvard University Press.

\noindent Vogeler, C.~S. (2017). Why do territorial farm animal welfare policies
differ? Explaining regulatory variations in the EU's multi-level system.
\textit{Journal of Comparative Policy Analysis}, \textit{19}(1), 18--34.

\noindent Wallenbeck, A., H\"{o}\"{o}g, Y., \& Gunnarsson, S. (2024). Pig welfare
policy in Europe: Adoption sequences, gold-plating, and member state divergence.
\textit{Animal}, \textit{18}(4), 101122.

\noindent World Animal Protection. (2020). \textit{Animal Protection Index 2020}.
World Animal Protection. \url{https://api.worldanimalprotection.org/}

\noindent Z\"{o}llmer, J. (2025). Race to the top of farm animal welfare policies in
US states. \textit{Journal of Comparative Policy Analysis}.
\url{https://doi.org/10.1080/13876988.2025.2493820}

\endgroup

\end{document}
